\section{Introduction et objectifs de la revue}

L'utilisation de préparations issues des feuilles d'olivier chez le poulet de chair se situe à l'interface de plusieurs domaines : physiologie digestive aviaire, nutrition et rationnement, santé intestinale, phytothérapie, valorisation des coproduits oléicoles et chimie des composés phénoliques. Une partie bibliographique cohérente doit donc présenter d'abord le contexte biologique et nutritionnel du poulet, puis la matière première végétale et ses procédés de transformation, avant d'examiner les effets réellement observés chez l'animal.

Le présent manuscrit adopte volontairement une stratégie de couverture large. Les thèmes retrouvés dans les thèses vétérinaires et documents de référence fournis sont intégrés lorsqu'ils peuvent contribuer à la compréhension du sujet : particularités anatomo-physiologiques digestives, microbiote, besoins alimentaires, rationnement, facteurs de croissance et leurs alternatives, importance de l'olivier en Tunisie, caractéristiques des feuilles, polyphénols, extraction et propriétés biologiques. Cette version constitue ainsi une base de travail destinée à être hiérarchisée et éventuellement allégée après relecture par l'encadrement.

L'élargissement du plan ne doit toutefois pas réduire l'exigence scientifique. Le terme \og extrait de feuilles d'olivier\fg{} recouvre des préparations différentes : infusion distribuée dans l'eau, extrait obtenu avec un solvant, poudre incorporée à l'aliment, fraction purifiée ou molécule isolée. Ces produits ne présentent ni la même composition ni la même exposition. Les travaux directement consacrés à l'administration dans l'eau de boisson constituent donc le niveau de comparaison le plus proche du sujet, tandis que les autres formes servent à compléter l'interprétation \citep{Jabri2017,Oke2017,Erener2023}.

La question centrale demeure la suivante : dans quelles conditions les préparations de feuilles d'olivier peuvent-elles modifier les performances et les indicateurs de santé du poulet de chair, et quelle part des résultats publiés est transposable à une administration dans l'eau de boisson ? L'analyse examine successivement la préparation et la caractérisation du produit, la croissance et l'efficacité alimentaire, les paramètres sanguins et immunitaires, puis l'intégrité intestinale, le microbiote et l'activité fermentaire.

La synthèse distingue les résultats expérimentaux rapportés, les mécanismes biologiques proposés et les hypothèses qui restent à vérifier. Elle ne présume ni les conditions d'élevage, ni les analyses réalisées, ni les résultats de Mohamed Hamida. Les thèses antérieures servent à construire le périmètre et à repérer les chapitres attendus ; leur contenu n'est pas repris comme substitut aux articles et références scientifiques vérifiables.
